\documentclass[10pt,a4paper]{amsart}
\usepackage{amsmath,amssymb,amsthm,graphicx}
% Retained minimal preamble for the Stacks Project chapter more-algebra.tex (More on Algebra).
% The chapter file pinned-original-source.tex carries only \input{preamble} on its first line, so no
% preamble text lives inside the chapter itself.  The Stacks Project uses one single shared preamble for
% every chapter; the snapshot stages it identically as sources/stacks-chow/preamble.tex and
% sources/stacks-etale-cohomology/preamble.tex.  The block below is copied mechanically, line for line,
% from sources/stacks-chow/preamble.tex lines 14-18 (the Xy-pic setup) and 154-217 (the native theorem
% declarations together with the mathematical macros).  Nothing was added, removed or redefined; the
% wrapper class and the amsmath/amssymb/amsthm/graphicx base packages are supplied by the assembler.
% Chapter: more-algebra.tex -- More on Algebra.
% For commutative diagrams we use Xy-pic
\usepackage[all]{xy}

% We use 2cell for 2-commutative diagrams.
\xyoption{2cell}

% Theorem environments.
%
\theoremstyle{plain}
\newtheorem{theorem}[subsection]{Theorem}
\newtheorem{proposition}[subsection]{Proposition}
\newtheorem{lemma}[subsection]{Lemma}

\theoremstyle{definition}
\newtheorem{definition}[subsection]{Definition}
\newtheorem{example}[subsection]{Example}
\newtheorem{exercise}[subsection]{Exercise}
\newtheorem{situation}[subsection]{Situation}

\theoremstyle{remark}
\newtheorem{remark}[subsection]{Remark}
\newtheorem{remarks}[subsection]{Remarks}

\numberwithin{equation}{subsection}

% Macros
%
\def\lim{\mathop{\mathrm{lim}}\nolimits}
\def\colim{\mathop{\mathrm{colim}}\nolimits}
\def\Spec{\mathop{\mathrm{Spec}}}
\def\Hom{\mathop{\mathrm{Hom}}\nolimits}
\def\Ext{\mathop{\mathrm{Ext}}\nolimits}
\def\SheafHom{\mathop{\mathcal{H}\!\mathit{om}}\nolimits}
\def\SheafExt{\mathop{\mathcal{E}\!\mathit{xt}}\nolimits}
\def\Sch{\mathit{Sch}}
\def\Mor{\mathop{\mathrm{Mor}}\nolimits}
\def\Ob{\mathop{\mathrm{Ob}}\nolimits}
\def\Sh{\mathop{\mathit{Sh}}\nolimits}
\def\NL{\mathop{N\!L}\nolimits}
\def\CH{\mathop{\mathrm{CH}}\nolimits}
\def\proetale{{pro\text{-}\acute{e}tale}}
\def\etale{{\acute{e}tale}}
\def\QCoh{\mathit{QCoh}}
\def\Ker{\mathop{\mathrm{Ker}}}
\def\Im{\mathop{\mathrm{Im}}}
\def\Coker{\mathop{\mathrm{Coker}}}
\def\Coim{\mathop{\mathrm{Coim}}}

% Boxtimes
%
\DeclareMathSymbol{\boxtimes}{\mathbin}{AMSa}{"02}

%
% Macros for moduli stacks/spaces
%
\def\QCohstack{\mathcal{QC}\!\mathit{oh}}
\def\Cohstack{\mathcal{C}\!\mathit{oh}}
\def\Spacesstack{\mathcal{S}\!\mathit{paces}}
\def\Quotfunctor{\mathrm{Quot}}
\def\Hilbfunctor{\mathrm{Hilb}}
\def\Curvesstack{\mathcal{C}\!\mathit{urves}}
\def\Polarizedstack{\mathcal{P}\!\mathit{olarized}}
\def\Complexesstack{\mathcal{C}\!\mathit{omplexes}}
% \Pic is the operator that assigns to X its picard group, usage \Pic(X)
% \Picardstack_{X/B} denotes the Picard stack of X over B
% \Picardfunctor_{X/B} denotes the Picard functor of X over B
\def\Pic{\mathop{\mathrm{Pic}}\nolimits}
\def\Picardstack{\mathcal{P}\!\mathit{ic}}
\def\Picardfunctor{\mathrm{Pic}}
\def\Deformationcategory{\mathcal{D}\!\mathit{ef}}

% --- appended 2026-10-12 (Chief): the project own definitions of the reference/slogan/history
% display environments, copied mechanically from the shared Stacks preamble lines 8-12
% (\usepackage{verbatim} plus the three \newenvironment lines, which the chapter preamble keeps
% verbatim from the project preamble). Nothing added, removed or redefined: these environments
% carry no mathematical content and the project itself typesets them as comments. Without these
% three lines a counted statement whose own text uses them fails with Environment slogan undefined.
\usepackage{verbatim}
\newenvironment{reference}{\comment}{\endcomment}
\newenvironment{slogan}{\comment}{\endcomment}
\newenvironment{history}{\comment}{\endcomment}
\usepackage{cleveref}
\crefname{originallemma}{original lemma}{original lemmas}
\Crefname{originallemma}{Original lemma}{Original lemmas}
\makeatletter
\newlabel{derived-definition-K-zero}{{0FCN}{source}{}{sourcepointer-1}{}}
\newlabel{derived-definition-K-zero@cref}{{[originallemma][0][]0FCN}{[1][0][]source}{}{}{}}
\makeatother
\crefname{originallemma}{original lemma}{original lemmas}
\Crefname{originallemma}{Original lemma}{Original lemmas}
\makeatletter
\newlabel{derived-section-canonical-delta-functor}{{014Z}{source}{}{sourcepointer-2}{}}
\newlabel{derived-section-canonical-delta-functor@cref}{{[originallemma][0][]014Z}{[1][0][]source}{}{}{}}
\makeatother
\crefname{originallemma}{original lemma}{original lemmas}
\Crefname{originallemma}{Original lemma}{Original lemmas}
\makeatletter
\newlabel{homology-section-truncations}{{0118}{source}{}{sourcepointer-3}{}}
\newlabel{homology-section-truncations@cref}{{[originallemma][0][]0118}{[1][0][]source}{}{}{}}
\makeatother
\crefname{originallemma}{original lemma}{original lemmas}
\Crefname{originallemma}{Original lemma}{Original lemmas}
\makeatletter
\newlabel{lemma-perfect-ring-classical-generator}{{20288}{source}{}{sourcepointer-4}{}}
\newlabel{lemma-perfect-ring-classical-generator@cref}{{[originallemma][0][]20288}{[1][0][]source}{}{}{}}
\makeatother
\crefname{originallemma}{original lemma}{original lemmas}
\Crefname{originallemma}{Original lemma}{Original lemmas}
\makeatletter
\newlabel{lemma-perfect-to-K-group}{{36491}{source}{}{sourcepointer-5}{}}
\newlabel{lemma-perfect-to-K-group@cref}{{[originallemma][0][]36491}{[1][0][]source}{}{}{}}
\makeatother
\title{The Euler characteristic identifies the K-zero of the perfect complexes of a ring with the K-zero of the ring}
\author{The Stacks Project Authors}
\date{stacks-project:more-algebra.tex:local-snapshot-20261012; GNU FDL 1.2 or later}
\begin{document}
\maketitle
\setcounter{subsection}{1}
\noindent\textit{Editorial note.} All delivered mathematical text is the author's own, copied from the pinned Stacks Project chapter file more-algebra.tex; no statement or proof step was rewritten, reordered or removed.\par\medskip
\noindent\textit{Editorial note.} The chapter file pinned-original-source.tex carries only \textbackslash{}input\{preamble\} on its first line, so no preamble text lives inside the chapter. The retained file sources/stacks-more-algebra/retained-preamble.tex is a mechanical, line-for-line copy of the project's single shared preamble as staged in the snapshot; nothing was added, removed or redefined and no package was added.\par\medskip
\noindent\textit{Editorial note.} The source URL is the source-level Stacks Project URL, as required when more than one card is produced from one source; the per-claim tag is carried by claim\_key.\par\medskip
\noindent\textit{Editorial note.} The author writes 'with apologies for the horrendous notation' and, in the next paragraph, 'then the the existence'; both are reproduced unchanged. The five references inside the counted proof span are recorded as non-load-bearing navigation pointers.\par\medskip
\noindent\textit{Allowed standard prerequisites (editorial).}\begin{itemize}
\item K\_0 of an additive or triangulated category by generators and relations. invoked at proof lines 36588-36590 and 36596-36598. Fixes the two groups being compared and makes both displayed maps well defined.
\item Distinguished triangles from short exact sequences of projective modules, and from truncations of a bounded complex. used at proof lines 36597-36598 and 36604-36619. Is the tool that lets a statement about modules be promoted to a statement about complexes.
\item The Euler characteristic map c of the preceding lemma and its additivity on distinguished triangles. invoked at statement line 36583. Supplies the map whose universal property is being established.
\end{itemize}\medskip
\ifdefined\hypertarget\hypertarget{sourcepointer-1}{}\fi
\noindent\textit{Source reference.} Original lemma 0FCN: focused-seeds/AGP/sources/stacks-snapshot-20261011/tags.txt, source lines 19188--19188. Cross-chapter navigation pointer into the Stacks Project; the standard fact it supplies is applied inside the counted proof span and the argument resting on it stays there.\par\medskip
\ifdefined\hypertarget\hypertarget{sourcepointer-2}{}\fi
\noindent\textit{Source reference.} Original lemma 014Z: focused-seeds/AGP/sources/stacks-snapshot-20261011/tags.txt, source lines 1490--1490. Cross-chapter navigation pointer into the Stacks Project; the standard fact it supplies is applied inside the counted proof span and the argument resting on it stays there.\par\medskip
\ifdefined\hypertarget\hypertarget{sourcepointer-3}{}\fi
\noindent\textit{Source reference.} Original lemma 0118: focused-seeds/AGP/sources/stacks-snapshot-20261011/tags.txt, source lines 1354--1354. Cross-chapter navigation pointer into the Stacks Project; the standard fact it supplies is applied inside the counted proof span and the argument resting on it stays there.\par\medskip
\ifdefined\hypertarget\hypertarget{sourcepointer-4}{}\fi
\noindent\textit{Source reference.} Original lemma 20288: sources/stacks-more-algebra/pinned-original-source.tex, source lines 20288--20288. Same-chapter navigation pointer; the statement it names is standard for the section and the argument that consumes it stays inside the counted proof span.\par\medskip
\ifdefined\hypertarget\hypertarget{sourcepointer-5}{}\fi
\noindent\textit{Source reference.} Original lemma 36491: sources/stacks-more-algebra/pinned-original-source.tex, source lines 36491--36491. Same-chapter navigation pointer; the statement it names is standard for the section and the argument that consumes it stays inside the counted proof span.\par\medskip
\begin{lemma}
\label{lemma-perfect-to-K-group-universal}
Let $R$ be a ring. Let $D_{perf}(R)$ be the derived category of
perfect objects, see Lemma \ref{lemma-perfect-ring-classical-generator}.
The map $c$ of Lemma \ref{lemma-perfect-to-K-group} gives an isomorphism
$K_0(D_{perf}(R)) = K_0(R)$.
\end{lemma}

\begin{proof}
It follows from the definition of $K_0(D_{perf}(R))$
(Derived Categories, Definition \ref{derived-definition-K-zero})
that $c$ induces a homomorphism $K_0(D_{perf}(R)) \to K_0(R)$.

\medskip\noindent
Given a finite projective module $M$ over $R$ let us denote
$M[0]$ the perfect complex over $R$ which has $M$ sitting in degree $0$
and zero in other degrees. Given a short exact sequence
$0 \to M \to M' \to M'' \to 0$ of finite projective modules
we obtain a distinguished triangle $M[0] \to M'[0] \to M''[0] \to M[1]$, see
Derived Categories, Section \ref{derived-section-canonical-delta-functor}.
This shows that we obtain a map $K_0(R) \to K_0(D_{perf}(R))$
by sending $[M]$ to $[M[0]]$ with apologies for the horrendous notation.

\medskip\noindent
It is clear that $K_0(R) \to K_0(D_{perf}(R)) \to K_0(R)$ is the identity.
On the other hand, if $M^\bullet$ is a bounded complex of finite projective
$R$-modules, then the the existence of the distinguished triangles
of ``stupid truncations''
(see Homology, Section \ref{homology-section-truncations})
$$
\sigma_{\geq n}M^\bullet \to \sigma_{\geq n - 1}M^\bullet \to
M^{n - 1}[-n + 1] \to (\sigma_{\geq n}M^\bullet)[1]
$$
and induction show that
$$
[M^\bullet] = \sum (-1)^i[M^i[0]]
$$
in $K_0(D_{perf}(R))$ (with again apologies for the notation).
Hence the map $K_0(R) \to K_0(D_{perf}(R))$ is surjective which
finishes the proof.
\end{proof}

\end{document}
